\scanpage{280}
\mspage{196}
% archive sheet number: 279
\[ up_{\alpha}=u_{\alpha}. \]

Or, en appliquant la propriété a) à $X_\alpha$ au lieu de $X$, on voit que la famille des flèches

\[ i(\deleted{\unsure{x}}y)\xrightarrow{\deleted{\unsure{p_y}}q_y}X_\alpha\ \unsure{,|\phi} \]

(indexée par $C_{/X_\alpha}$) est épim. stricte, a fortiori épimorphique. Donc pour prouver $up_\alpha=u_\alpha$, il suffit de prouver qu'on a

{\makeatletter\tagsleft@true\makeatother
\[ (up_\alpha)q_y=u_\alpha q_y \tag*{$(**)$} \]
}

pour \unsure{tt|tout} $y$\deleted{\unsure{$\in\mathrm{Ob}$}} dans $C_{/X_\alpha}$. Mais c'est là une conséquence de la compatibilité ($\alpha$) (\unsure{deuxième|dernière} \unsure{formule}), appliquée au cas suivant\unsure{:}

\diagram{\begin{tikzcd}[ampersand replacement=\&,row sep=small]
\& i(y) \arrow[dr,dashed,"p_\alpha q_y=p_x"] \&\\
Z=i(y) \arrow[ur,"\mathrm{id}"] \arrow[dr,"q_y"'] \& \& X{,}\\
\& X_\alpha \arrow[ur,dashed,"p_\alpha"'] \&
\end{tikzcd}}

on a donc, en posant $\deleted{\unsure{x_\alpha}}x=(y,p_\alpha q_y)\in\mathrm{Ob}\,C_{/X}$,

\[ \underbrace{u_{\unsure{y|x}}\circ\mathrm{id}_{i(y)}}_{u_x}=u_\alpha q_y, \]

donc pour prouver $\unsure{(**)|(\alpha)}$, il reste à vérifier

\[ u_x=\underbrace{(u\deleted{\unsure{q}}p_\alpha)q_y)}_{u(p_\alpha q_y)=up_x}\deleted{\unsure{\mathrm{id}}}, \]
